\section{Business Understanding}

\subsection{Determine Business Objectives}

Modern code assistants and automated review systems require data that describes how software changes are actually produced. Public code datasets usually contain source files, commits, or code--text pairs, but they often miss the surrounding workflow: the pull request, review discussion, requested changes, and final diff. This project addresses that gap by constructing a curated dataset of public GitHub pull request workflows.

\textbf{Business objective:} create a reusable dataset product for machine learning teams that build code review automation, coding assistants, and developer-agent systems. Each record should represent one pull request workflow and combine pull request metadata, review activity, discussion, changed files, code diff, quality metadata, and provenance.

\textbf{Business success criteria:} the project is successful if it produces a multi-repository dataset with at least 5,000 high-quality workflows, removes obvious noise such as bots and generated-only changes, supports a realistic downstream prediction task, and demonstrates that the dataset contains learnable signal beyond simple baselines.

\subsection{Assess Situation}

The main data source is GH Archive, which provides public GitHub event history and makes it possible to discover merged pull request candidates at scale. The GitHub REST API is then used to enrich each candidate with complete pull request metadata, changed files, reviews, discussion comments, and unified diffs. The work is constrained to public repositories and local compute, so the collection strategy must be efficient and reproducible.

The main risk is data noise. Public GitHub activity contains bots, trivial edits, generated files, documentation-only changes, and workflows with too little review content. This risk is handled by filtering, quality metadata, and repository-level splitting. Another important risk is label imbalance: the later modeling target, \texttt{review\_concern}, is positive for most records. Therefore, accuracy alone is not a meaningful success measure; macro-F1, balanced accuracy, and ROC-AUC are more informative.

The expected value of the dataset is saved engineering effort. A team that wants workflow-level GitHub data would otherwise need to implement event discovery, API enrichment, filtering, schema design, and quality checks. A compact but curated dataset of roughly 10,000 enriched workflows can therefore be positioned as a niche dataset product for AI engineering teams.

\subsection{Determine Data Mining Goals}

The first data mining goal is corpus construction. The project should produce at least 5,000 accepted pull request workflows, with a target close to 10,000 raw enriched records. A workflow is accepted when it represents a merged pull request with source-code changes, non-empty source patches, review and discussion activity, and no strong bot, documentation-only, generated-only, or trivial-change pattern.

The second data mining goal is supervised prediction of \texttt{review\_concern}. The model receives only pull-request-intrinsic information available at opening time, including metadata, changed-file summaries, and capped code-diff embeddings. It predicts whether later review activity will contain a substantive concern signal, defined by a \texttt{CHANGES\_REQUESTED} review or at least two meaningful review comments.

\textbf{Data mining success criteria:} the prepared dataset should pass the size and quality thresholds, and the best held-out model should beat a majority-class baseline on macro-F1 and random ranking on ROC-AUC. These criteria are intentionally modest because the modeling stage is used to validate dataset usefulness, not to claim production readiness.

\subsection{Produce Project Plan}

The project follows the CRISP-DM sequence. First, the business objective and downstream task are fixed. Second, candidate pull requests are collected and enriched. Third, the data is explored for scale, completeness, size, review activity, language diversity, and quality limitations. Fourth, the raw records are converted into a modeling table with leakage-safe features and repository-level splits. Fifth, several baseline and non-linear models are trained and compared. Finally, the results are evaluated against the technical success criteria and used to decide whether another iteration is required.

The selected tools are standard and local: Python for data processing, GH Archive and the GitHub REST API for data acquisition, Pandas and NumPy for data preparation, ModernBERT embeddings for compact diff representation, scikit-learn for baselines, and XGBoost for the main non-linear model.
